\scanpage{62}
Let us already set
\[
(14)\qquad\left\{\begin{aligned}
V_{\Delta_1}&=\coprod_{\substack{(i,j)\in I^2\\i\leq j}}V_{i,j}\\
\overset{\deleted{V}}{V}_{\Delta_1}^*&=\coprod_{\substack{(i,j)\in I^2\\i\leq j}}V_{i,j}^*\underset{\text{open}}{\subset}V_{\Delta_1}
\end{aligned}\right.
\]
we therefore have
\diagram{(15)\qquad\begin{tikzcd}[column sep=large,row sep=large]
V_{\Delta_1}^*\arrow[r,hook]\arrow[d] & V_{\Delta_1}\arrow[d,dashed] & X_{\Delta_0}^*\arrow[l]\arrow[dl,dashed]\\
\deleted{V}X_{\Delta_0}^*\arrow[r,dashed] & X &
\arrow[from=2-1,to=2-2,phantom,"("{pos=.12}]\arrow[from=2-2,to=2-3,phantom,")"{pos=.05}]
\end{tikzcd}}
\diagram{\deleted{\unsure{(16)}\qquad\begin{tikzcd}
V_{\Delta_1} & X_{\Delta_0^{\unsure{2}}}^*\arrow[l]
\end{tikzcd}}}
which reflect the diagrams
% The upper left arrow in (17) has a narrow double-looking shaft.
% Single versus doubled shaft is marked explicitly on the arrow.
\diagram{(17)\qquad\begin{tikzpicture}
\node[inner sep=5mm] (d) {\begin{tikzcd}[column sep=large,row sep=large]
V_{i,j}^*\arrow[r,phantom,"\unsure{\longrightarrow|\Longrightarrow}"{description}]\arrow[d] & V_{i,j}\arrow[d,dashed] & X_i^*\arrow[l]\arrow[dl,dashed]\\
X_j^*\arrow[r,dashed] & X & {.}
\arrow[from=2-1,to=2-2,phantom,"("{pos=.12}]\arrow[from=2-2,to=2-3,phantom,")"{pos=.05}]
\end{tikzcd}};
\draw ([xshift=1mm,yshift=1mm]d.north west)--([xshift=-2mm,yshift=2mm]d.north east);
\draw[->] ([yshift=-1mm]d.north west)--([xshift=.5mm]d.west);
\draw ([xshift=.5mm]d.west)--([xshift=1mm,yshift=1mm]d.south west);
\draw ([yshift=-1mm]d.north east)--([yshift=1mm]d.south east);
\draw ([xshift=3mm]d.south west)--([xshift=-1mm,yshift=1mm]d.south east);
\end{tikzpicture}}

\margin{\uline{NB} The form
\diagram{\begin{tikzcd}[column sep=small,row sep=small]
&\bullet\arrow[d]\\
\bullet\arrow[r]\arrow[d]&\bullet\\
\bullet&
\end{tikzcd}}
is perhaps more suggestive}

\deleted{\unsure{Let us note|We have}} In the end, one is not sure that
the notations (14) \unsure{are|are} really useful —
in any case the square in \unsure{(15)|(17)} has
nothing cocartesian about it, for the diagram in
solid lines, which is the ``sum'' diagram
\margin{61}
